% ============================================================
% DRAFT SNIPPET — state-vs-trait functional characterization
% Discussion subsection + appendix section for paper-v2-narrow.
% RESID numbers marked XX pending the axis-of-record rerun.
% ============================================================

% ---------- Discussion subsection (insert after "What the component
% ---------- result does and does not localize")

\subsection{The working axis behaves as transient state, not a persistent trait}
\label{sec:state-trait}

If the intervention axis is a behaviorally relevant mixture rather than a
construct-pure representation, one can still ask what \emph{functional role}
the mixture plays. Recent work distinguishes two roles for residual-stream
directions: locally scoped emotion-concept inputs to next-token prediction
\cite{sofroniew2026emotion}, and persona directions --- persistent latents
whose training-data projection shift predicts durable trait acquisition after
finetuning \cite{chen2025persona}. Because both families are difference-in-means
directions extracted by near-identical pipelines, the distinction is not
recoverable from the vector itself; it is a dynamical and causal property. We
test it with three null-calibrated probes (Appendix~\ref{app:state-trait}).

The working axis reads state-like on every probe. Its per-token projection
autocorrelation timescale on long generated texts is $\lambda = 48.5$ tokens
against a 50-random-direction null (median 61, 95th percentile 209) --- random
directions inherit slow residual-stream drift, so falling at or below the null
median indicates faster-than-drift, locally scoped content. A speaker-framing
contrast (Assistant-indexed versus third-person renderings of matched
scenarios) lands at the 32nd percentile of the same random-direction null;
surface-format differences alone produce large framing effects for arbitrary
directions, so only exceedance of that null would count as Assistant-indexing.
Finally, a finetuning test in the protocol of \cite{chen2025persona}: three
LoRA datasets with graded empathic content over identical trait-neutral
prompts produce a cleanly graded projection shift $\Delta P$
($+5.97$ / $+1.70$ / $-1.85$ along the axis for empathic / mixed /
non-empathic data), yet the post-finetuning \emph{chronic}
projection on affect-free text is not proportional to $\Delta P$
($\beta = 0.060$, 22nd percentile of the per-direction null; genuine persona
directions reach $r = 0.76$--$0.97$ on this protocol). A blinded behavioral
judge shows the same dissociation: finetuned models fully express the trained
disposition on in-distribution scenarios (composite $1.41 \to 3.00$ empathic,
$\to 0.57$ non-empathic) while generic out-of-distribution prompts move
within noise ($2.00 \to 2.02$).

The allowed reading is narrow: the axis carries transiently scoped,
input-driven content, and finetuning on empathy-graded data at this scale
(LoRA $r{=}16$, 400 examples, 2 epochs) does not install a chronic setpoint
along it. We do not claim the axis \emph{is} an emotion representation, nor
that heavier finetuning could not convert its role.

% ---------- Appendix section

\section{State-Versus-Trait Functional Battery}
\label{app:state-trait}

All three probes share one null construction: every statistic computed for
the candidate axis is computed identically for norm-matched random directions
run through the same inputs, and verdicts are percentiles in that null. This
matters because both alternative absolute criteria fail: raw autocorrelation
timescales are dominated by slow residual-stream drift (random-direction
median up to $\sim$100 tokens), and Assistant-versus-third-person framings
differ in surface format (dialogue markup, length, register), giving random
directions framing effects up to $|d| \approx 3$.

\paragraph{Timescale.}
Per-token projections $z_t = a_t \cdot \hat{v}$ over 60 generated texts
(median 294 words, both contrast polarities); exponential fit
$\rho(\tau) \approx e^{-\tau/\lambda}$, $\tau \le 64$. Candidate $\lambda$
is compared with 50 random directions on the same texts.

\paragraph{Speaker invariance.}
Sixteen scenario items rendered both as third-person narration (four subject
variants) and as Assistant dialogue turns with matched concept content;
Cohen's $d$ of last-token projections between framings, candidate versus the
same 50-direction null.

\paragraph{Finetuning ($\beta$) test.}
Three LoRA datasets (400 examples each, balanced over the five scenario
families) hold prompts fixed --- scenario and objective only, with the
generation-time trait instruction removed --- and vary only responses:
empathic, non-empathic, or an equal mixture. $\Delta P$ is the mean
response-token projection of dataset responses minus the base model's own
responses to the same prompts; the base-response term is constant across
datasets and cancels in the slope. Two seeds per dataset
(losses declined in all six runs). The durable-shift readout is the mean
token projection on 30 affect-free neutral prompts;
$\beta$ is the least-squares slope of that shift against $\Delta P$ over the
three datasets plus the un-finetuned anchor at the origin. The specificity
null computes the same $\beta$ for 40 random directions, each along itself,
on the same checkpoints. A blinded judge (Claude Haiku batch, temperature 0,
checkpoint identities hidden behind shuffled opaque ids; 315/315 scored)
rates every checkpoint generation on the repository's standard 0--4
behavioral dimensions.

\begin{table}[h]
\centering
\caption{State-versus-trait battery. Percentiles are within the
random-direction null for the same statistic; exceeding the null 95th
percentile would indicate trait-like persistence (timescale), Assistant
indexing (invariance), or $\Delta P$-predicted durable shift ($\beta$).}
\label{tab:state-trait}
\small
\begin{tabular}{lcccc}
\toprule
Direction & $\lambda$ (null med / p95) & inv.\ pct & $\beta$ (null med / p95) & $\beta$ pct \\
\midrule
DFA L8            & 46.4 (102 / 582)  & 90 & 0.028 (0.167 / 1.442) & 15 \\
DFA L20           & 52.7 (61 / 209)   & 14 & 0.030 (0.167 / 1.442) & 18 \\
M block 20        & 7.7 (61 / 209)    & 40 & 0.077 (0.167 / 1.442) & 30 \\
Residualized axis & 48.5 (61 / 209)   & 32 & 0.060 (0.167 / 1.442) & 22 \\
\bottomrule
\end{tabular}
\end{table}

\begin{table}[h]
\centering
\caption{Blinded behavioral judgments (composite of helping action,
affective language, caring character; 0--4 scale). Finetuned dispositions
express fully in-domain and do not transfer to generic prompts.}
\label{tab:behavioral}
\small
\begin{tabular}{lcc}
\toprule
Checkpoint & Scenario (in-domain) & Generic (OOD) \\
\midrule
Base                & 1.41 & 2.00 \\
Empathic (2 seeds)  & 3.01 / 2.99 & 2.12 / 2.02 \\
Mixed (2 seeds)     & 1.89 / 1.57 & 2.03 / 2.02 \\
Non-empathic (2 seeds) & 0.56 / 0.57 & 1.63 / 1.65 \\
\bottomrule
\end{tabular}
\end{table}

% ---------- references.bib additions
% @article{chen2025persona, title={Persona Vectors: Monitoring and
%   Controlling Character Traits in Language Models}, author={Chen, Runjin
%   and others}, journal={arXiv preprint arXiv:2507.21509}, year={2025}}
% @article{sofroniew2026emotion, title={Emotion Concepts and their Function
%   in a Large Language Model}, author={Sofroniew, Nicholas and Kauvar,
%   Isaac and others}, journal={arXiv preprint arXiv:2604.07729}, year={2026}}
